\emph{Learned continuous communication.} CommNet~\cite{sukhbaatar2016learning} and DIAL~\cite{foerster2016learning} train messages end to end by backpropagating through the channel; TarMAC~\cite{das2018tarmac} adds targeted attention. Our continuous baseline follows the DIAL idea (differentiable, noisy channel) but is a simplified reimplementation, not the published system. \emph{Discrete and emergent language.} Lazaridou et al.~\cite{lazaridou2016multi}, Havrylov and Titov~\cite{havrylov2017emergence} and Mordatch and Abbeel~\cite{mordatch2017emergence} study symbols that emerge in cooperative games; discrete channels are typically trained with a continuous relaxation such as Gumbel-softmax~\cite{jang2016categorical}, as in our discrete baseline. Kottur et al.~\cite{kottur2017natural} show that emergent languages are not automatically natural or compositional, which motivates comparing against a fixed human-designed language. Lowe et al.~\cite{lowe2019pitfalls} caution that task success alone is a coarse measure of emergent communication; our no-communication baseline and cross-play control are the only checks here that messages are causally used and partner-specific, and we do not measure message semantics. \emph{Partner robustness.} Other-Play~\cite{hu2020other} shows that conventions learned in self-play can fail with independently trained partners, and trains agents to be robust to symmetry-equivalent relabellings of the problem; our cross-play protocol measures the failure it addresses, but we do not implement Other-Play. Multi-agent actor-critic~\cite{lowe2017multi} is a common reinforcement-learning framework for such tasks; we instead use pathwise gradients, a simplification discussed in the limitations.
